%% Chapitre 11 — Subagents et orchestration de tâches
%% RÈGLE : uniquement des commandes sémantiques bookforge. Aucun \chapter,
%% \textbf, \begin{lstlisting}, \vspace... brut.

\chapitre{Subagents et orchestration de tâches}

\introChapitre{
  Pour les gros chantiers, un seul agent dans une seule session ne suffit plus.
  Ce chapitre durable introduit les subagents et l'idée d'orchestration :
  déléguer des sous-parties à des agents spécialisés, garder une vue d'ensemble,
  et faire converger leurs résultats. C'est le saut du « j'utilise un agent » au
  « je pilote plusieurs agents ».
}

\section{Pourquoi un seul agent ne suffit pas}

\paragraphe{Tu sais maintenant découper un chantier en tâches délégables et
vérifier ce que l'agent produit. Tant que la tâche tient dans une session,
ça marche. Mais arrive un moment où elle déborde.}

\paragraphe{Le premier mur est le contexte. Un agent ne raisonne que sur ce
qu'il a sous les yeux, et sa fenêtre de contexte est finie. Sur un gros
chantier, à force de lire des fichiers et d'accumuler des résultats de
commandes, il sature. Les détails du début s'effacent, la qualité chute. Tu
le vois quand l'agent oublie une décision prise vingt minutes plus tôt.}

\paragraphe{Le deuxième mur est l'attention. Même sans saturer, un agent qui
jongle entre la migration de la base, le rafraîchissement de l'interface et la
mise à jour des tests fait tout moyennement. Il mélange les registres, perd le
fil d'un sujet en traitant l'autre. Comme un humain multitâche, il devient
médiocre partout.}

\paragraphe{Le troisième mur est le séquentiel. Dix modules indépendants à
réécrire, un agent les traite l'un après l'autre. Tu attends. Pourtant, rien
n'oblige à ce que ces dix chantiers se suivent : ils pourraient avancer de
front.}

\paragraphe{La réponse à ces trois murs est la même : \important{ne plus
mettre un seul agent sur tout}. On répartit. C'est exactement le geste que tu
fais déjà entre humains quand un projet dépasse une personne.}

\section{Subagents : déléguer une sous-tâche à un agent dédié}

\paragraphe{Un \terme{subagent} est un agent secondaire à qui l'agent
principal confie une sous-tâche délimitée, avec son \emphase{propre} contexte.
C'est le point clé : le subagent démarre frais, ne voit que ce qu'on lui donne,
travaille dans son coin, et renvoie un résultat.}

\paragraphe{Cette isolation du contexte est un atout, pas une contrainte.
Le subagent qui analyse la couverture de tests n'encombre pas la session
principale avec les centaines de lignes qu'il a lues : il remonte une synthèse.
L'agent principal récupère la conclusion, pas le fouillis. Sa propre fenêtre
de contexte reste légère et nette.}

\paragraphe{Le revers : le subagent ne sait que ce que tu lui transmets dans
son prompt de lancement. Il n'a pas vécu la conversation. Une sous-tâche mal
cadrée — but flou, contexte manquant, critère de succès absent — échoue plus
durement encore qu'avec un agent ordinaire, parce qu'il n'y a personne pour
rattraper en cours de route. La règle de la tâche délégable du chapitre~04
s'applique ici à la lettre.}

\begin{astuce}
Délègue à un subagent ce qui est \emphase{borné} et dont seul le résultat
compte : « audite l'usage de cette API dépréciée et liste les fichiers à
changer », « écris les tests de ce module ». Garde dans la session principale
ce qui demande le fil de la conversation et tes arbitrages.
\end{astuce}

\section{Spécialiser des agents par rôle ou par domaine}

\paragraphe{Le contexte d'un subagent ne contient pas que la tâche : il
contient aussi un \emphase{rôle}. Tu peux préfixer son prompt par une posture —
« tu es un relecteur de sécurité », « tu es un rédacteur de documentation » —
et lui donner des consignes propres à ce rôle. Le même modèle, cadré
différemment, produit un travail différent.}

\paragraphe{Deux axes de spécialisation reviennent. Par \important{rôle} :
un agent qui écrit, un autre qui relit, un troisième qui teste. Le relecteur
n'a pas le biais de l'auteur ; il attaque le code avec un regard neuf, ce qui
le rend bien plus efficace pour traquer une hallucination glissée dans la
première passe.}

\paragraphe{Par \important{domaine} : un agent pour le \etranger{back-end},
un pour le \etranger{front-end}, un pour les migrations de données. Chacun
reçoit en contexte les conventions et les fichiers de son périmètre, et rien
d'autre. Plus le périmètre est étroit, plus le contexte est pertinent, et
meilleur est le résultat.}

\paragraphe{Quand tu retrouves toujours les mêmes rôles, fige-les. Un rôle de
subagent réutilisable n'est rien d'autre qu'un prompt soigné — la même logique
qu'une slash command, mais pour une posture entière plutôt qu'une commande.
Tu l'écris une fois, tu le rejoues partout.}

\section{Le rôle d'orchestrateur : découper, distribuer, recoller}

\paragraphe{Dès qu'il y a plusieurs subagents, il faut quelqu'un au-dessus
pour les coordonner. C'est l'\terme{orchestration}, et l'agent qui s'en charge
est l'orchestrateur. Son travail tient en trois temps.}

\paragraphe{\important{Découper.} L'orchestrateur transforme le chantier en
sous-tâches autonomes. La qualité du découpage décide de tout : des morceaux
qui se chevauchent produiront des conflits ; des morceaux qui dépendent les uns
des autres ne pourront pas avancer en parallèle. Le bon découpage minimise les
liens entre les parties.}

\paragraphe{\important{Distribuer.} Il lance chaque subagent avec un prompt
autoporteur — but, contexte, périmètre, critère de succès — et le bon rôle.
L'orchestrateur ne fait pas le travail : il le confie et le suit.}

\paragraphe{\important{Recoller.} C'est l'étape qu'on sous-estime. Chaque
subagent renvoie un fragment ; il faut les assembler en un tout cohérent,
détecter les contradictions, combler les trous. Le recollage est aussi le
moment de la vérification : l'orchestrateur relit les diffs remontés avant de
les intégrer, exactement comme tu relis ceux d'un agent unique.}

\begin{avertissement}
L'orchestrateur hérite des défauts qu'il coordonne. Si un subagent hallucine
et que personne ne relit le fragment qu'il remonte, l'erreur entre dans
l'assemblage final, noyée parmi les autres. Plus il y a d'agents, plus la
vérification au recollage est indispensable, pas optionnelle.
\end{avertissement}

\section{Travail en parallèle : isolation et synchronisation}

\paragraphe{Le parallélisme est le gain le plus visible : dix subagents qui
travaillent en même temps finissent en une fraction du temps d'un agent seul.
Mais le parallélisme n'est sûr qu'à une condition : l'\important{isolation}.
Deux agents qui modifient le même fichier en même temps se marchent dessus.}

\paragraphe{La parade est de donner à chacun un terrain à lui. En pratique :
un dossier distinct, une copie de travail séparée, une branche dédiée — bref,
un garde-fou d'isolation au sens du chapitre~07, mais multiplié. Si les
périmètres ne se recouvrent pas, le parallélisme ne provoque aucun conflit,
par construction.}

\paragraphe{Reste la \important{synchronisation} : les points où les fils
parallèles doivent se rejoindre. Tu ne peux pas assembler tant que tous les
subagents n'ont pas fini ; tu ne peux pas tester l'ensemble avant l'assemblage.
Ces points d'attente sont normaux. L'art de l'orchestration consiste à en
avoir le moins possible — c'est-à-dire à découper de sorte que les parties
restent indépendantes le plus longtemps.}

\paragraphe{Et quand une dépendance est inévitable, on l'assume : telle partie
attend telle autre. Vouloir tout paralléliser à n'importe quel prix recrée les
conflits qu'on cherchait à fuir.}

\section{Cas vécu : un chantier mené avec plusieurs agents}

\paragraphe{Prenons un cas concret et méta : produire un livre technique d'une
quinzaine de chapitres avec des agents. Un seul agent aurait saturé son
contexte bien avant le tiers, et chaque chapitre aurait pâti de ce qui le
précédait dans la session.}

\paragraphe{Le découpage tombe sous le sens : un chapitre est une sous-tâche
quasi autonome. L'orchestrateur prépare donc un contrat partagé — ton,
vocabulaire imposé, conventions — puis lance \emphase{un subagent rédacteur
par chapitre}, en parallèle. Chaque subagent reçoit le plan de \emphase{son}
chapitre, le contrat commun, et écrit \emphase{uniquement} dans son dossier.}

\paragraphe{L'isolation est totale : un dossier par chapitre, donc aucun
conflit d'écriture, même avec quinze agents simultanés. Les subagents ne se
parlent pas — et n'ont pas à le faire. Leur cohérence ne vient pas d'une
communication, mais du contrat partagé que chacun a reçu en entier. C'est ce
qui remplace la conversation commune qu'ils n'ont jamais eue.}

\paragraphe{Le point de synchronisation arrive à la fin : quand tous les
chapitres sont rédigés, l'orchestrateur recolle. Il vérifie qu'aucun terme
n'est défini deux fois, que les transitions entre chapitres voisins
s'enchaînent, qu'un chapitre ne contredit pas le suivant. Là où le contrat
n'a pas suffi, il corrige.}

\paragraphe{Le bilan est instructif. Le parallélisme a divisé le temps total
par un grand facteur. La spécialisation a donné des chapitres mieux ciblés
qu'un agent fatigué n'en aurait écrit. Mais le travail de l'orchestrateur — le
contrat en amont, le recollage en aval — a pesé plus lourd que prévu. C'est la
leçon générale : \important{orchestrer ne supprime pas le travail, il le
déplace} vers le découpage et l'assemblage.}

\resumeChapitre{
  \element{\important{Subagent} : un agent secondaire avec son propre contexte vierge, à qui on confie une sous-tâche délimitée — il remonte un résultat, pas un fouillis.}
  \element{\important{Prompt autoporteur} : le subagent n'a pas vécu la conversation principale ; son brief doit contenir but, contexte, périmètre et critère de succès sans rien supposer.}
  \element{\important{Orchestrateur} : découpe le chantier, distribue les sous-tâches, puis recolle les fragments en vérifiant chaque diff remontés avant intégration.}
  \element{\important{Isolation pour le parallélisme} : donne à chaque subagent un terrain à lui (dossier, branche) — deux agents sur le même fichier se marchent dessus.}
  \element{\important{Orchestrer déplace le travail} : le gain en parallélisme se paie en travail de découpage en amont et de recollage en aval — anticipe-le, ne le sous-estime pas.}
}

\paragraphe{Tu sais désormais piloter plusieurs agents : les spécialiser,
les lancer en parallèle, les isoler, recoller leurs résultats. C'est puissant,
et c'est précisément quand on tient un outil puissant qu'on se blesse. Le
chapitre suivant rassemble les pièges et les cas vécus — les ratés assumés —
pour que les tiens coûtent moins cher.}
